% !TEX program = lualatex
% Auto-generated from syllabus — 08: 情報の加工1：文書・スライド

\documentclass[aspectratio=43,professionalfonts,handout,10pt]{beamer}
\usepackage{beamer_template}

\newcommand{\LectureNum}{08}
\newcommand{\LectureTitle}{情報の加工1：文書・スライド}
\newcommand{\CourseName}{情報リテラシー}
\newcommand{\TextPages}{-}

\begin{document}
% --- Slide 1: title ---

\begin{frame}{\CourseName\quad 第\LectureNum 回「\LectureTitle 」}
  {\small \TermName\quad \TeacherName\quad ／\quad 教科書 \TextPages}

  \vfill

  \begin{block}{今日の流れ}
    \begin{enumerate}
      \item ローカル保存とクラウド保存の違い
      \item PowerPointの基本操作と保存設定
      \item Word・PowerPointで自己紹介を作成・提出
    \end{enumerate}
  \end{block}
  \vfill

  \begin{block}{今日の目標}
    \begin{itemize}
      \item オフィスソフトの基本的な使い方が分かる。
      \item オフィスソフトを用いて文書の作成や発表資料の作成を行うことができる。
    \end{itemize}
  \end{block}
  \vfill

  \begin{exampleblock}{キーワード}
    \small \textbf{ローカル保存} ／ \textbf{クラウド保存} ／ \textbf{Word} ／ \textbf{PowerPoint} ／ \textbf{OneDrive}
  \end{exampleblock}
\end{frame}

% --- Slide 2: ローカル保存 vs クラウド保存 ---

\begin{frame}{ローカル保存 vs クラウド保存}
  \begin{columns}[T]
    \begin{column}{0.48\textwidth}
      \begin{block}{ローカル保存（PC内）}
        \small
        \begin{itemize}
          \item 高速・安定・通信不要
          \item ネット環境に左右されない
          \item 他のPCからはアクセス不可
          \item 高専では：基本はローカル保存
        \end{itemize}
      \end{block}
    \end{column}
    \begin{column}{0.48\textwidth}
      \begin{exampleblock}{クラウド保存（OneDrive等）}
        \small
        \begin{itemize}
          \item どこからでもアクセス可能
          \item 常時通信・通信速度に依存
          \item みんなで使うと学内ネットが重くなる
          \item 高専では：バックアップとして活用
        \end{itemize}
      \end{exampleblock}
    \end{column}
  \end{columns}

  \begin{alertblock}{演習1}
    ローカル保存とクラウド保存の違いを説明せよ
  \end{alertblock}

  \begin{slidenote}
    ローカル保存＝高速・安定だがPC限定。クラウド保存＝どこからでもアクセスできるが通信に依存。高専では基本ローカル保存，OneDriveはバックアップ用途。全員がクラウド保存すると学内ネットが重くなるのが実務上の理由。
  \end{slidenote}
\end{frame}

% --- Slide 3: OneDriveの使い方 ---

\begin{frame}{OneDriveの役割と使い方}
  \begin{block}{OneDrive＝定期的なバックアップ用}
    \begin{itemize}
      \item 普段の作業は\kw{ローカル保存}で行う
      \item OneDriveには\kw{定期的にバックアップ}としてコピーする
      \item 常時保存先としては使わない（通信負荷・速度低下の原因になる）
    \end{itemize}
  \end{block}

  \vfill

  \begin{noticebox}[注意]
    OneDriveを常時保存先にすると，ファイルを開くたびに通信が発生し，学内ネットワークが遅くなる原因になる。
  \end{noticebox}

  \begin{slidenote}
    OneDriveの「常時同期」と「バックアップとしてのコピー」の違いを理解させる。常時同期はファイル操作のたびに通信が走るため，多人数が同時に使うと学内ネットが逼迫する。週1回程度，重要ファイルをコピーする運用が望ましい。
  \end{slidenote}
\end{frame}

% --- Slide 4: PowerPointの基本操作 ---

\begin{frame}{PowerPointの基本操作と保存設定}
  {\small PowerPointを起動し，保存設定の変更と基本操作を体験する。}

  \vfill

  \begin{block}{手順}
    \begin{enumerate}
      \item スタートメニューから「PowerPoint」を起動する
      \item 「新しいプレゼンテーション」を作成する
      \item 「ファイル」→「その他」→「オプション」→「保存」タブを開く
      \item 「既定でコンピューターに保存する」にチェック→「OK」
      \item タイトルスライドのタイトル欄に自分の名前を入力してみる
    \end{enumerate}
  \end{block}

  \begin{slidenote}
    「既定でコンピューターに保存する」を有効にすると，保存先がOneDriveではなくローカルのドキュメントフォルダになる。これにより学内ネットワークへの負荷を軽減できる。
  \end{slidenote}
\end{frame}

% --- Slide 4: Wordの保存設定と自己紹介文 ---

\begin{frame}{Wordのローカル保存設定}
  {\small Wordのローカル保存設定を行う。}

  \vfill

  \begin{block}{手順}
    \begin{enumerate}
      \item Wordを起動し，「ファイル」→「オプション」→「保存」タブを開く
      \item 「既定でコンピューターに保存する」にチェック→「OK」
    \end{enumerate}
  \end{block}

  \begin{slidenote}
    Wordも同様に「既定でコンピューターに保存する」設定を行う。PowerPointと同じ手順なので，自力でできるか確認させる。
  \end{slidenote}
\end{frame}

% --- Slide 5: ファイル管理の注意 ---

\begin{frame}{ファイル管理の注意事項}
  \begin{noticebox}[ファイル管理のルール]
    \begin{itemize}
      \item ファイル名には自分の名前を必ず入れる（誰のファイルか分かるように）
      \item 保存場所を必ず確認する（どこに保存したか覚えておく）
      \item 提出前にファイルが開けるか必ず確認する
      \item 期限を守る（Teamsの課題タブで期限を確認する習慣をつける）
    \end{itemize}
  \end{noticebox}

  \begin{slidenote}
    ファイル名に名前を入れる理由：提出後に教員側で誰のファイルか識別するため。保存場所の確認も習慣化させる。
  \end{slidenote}
\end{frame}

% --- Slide 7: 演習 ---

\begin{frame}{演習：自己紹介を作ろう}
  \begin{alertblock}{演習2}
    Wordで自己紹介文を作成せよ（名前・学年・出身地・趣味・好きな食べ物・将来の目標）。ファイル名は「自己紹介\_自分の名前.docx」とすること。
  \end{alertblock}

  \vfill

  \begin{alertblock}{演習3}
    PowerPointで自己紹介スライドを1枚作成せよ。ファイル名は「自己紹介\_自分の名前.pptx」とすること。
  \end{alertblock}

  \begin{slidenote}
    Word＝詳しい説明・文章向き，PowerPoint＝視覚的・発表向きという用途の違いに気づかせる。両方作ることで，同じ内容でもソフトによって表現が変わることを体験させる。
  \end{slidenote}
\end{frame}

% --- チェックリスト ---

\begin{frame}{まとめ・チェックリスト}
  \begin{block}{自己チェック（できたら $\checkmark$ をつけよう）}
    \begin{itemize}
      \item[$\square$] ローカル保存とクラウド保存の違いを理解した
      \item[$\square$] PowerPointの保存設定を変更できた
      \item[$\square$] Wordで自己紹介文を作成・保存できた
      \item[$\square$] PowerPointで自己紹介スライドを作成・保存できた
      \item[$\square$] Teamsの課題タブから2つのファイルを提出できた
    \end{itemize}
  \end{block}

  \vfill

  \begin{alertblock}{演習4}
    演習2（Word）と演習3（PowerPoint）の2つのファイルをTeamsの課題タブから提出せよ
  \end{alertblock}
\end{frame}

\end{document}
